% ==============================================================================
% Fichier     : Cours.tex
% Titre       : Audit des Systèmes d'Information
% Auteur      : MINKA MI NGUIDJOÏ Thierry Emmanuel
% Institution : Université de Yaoundé I - Ecole Nationale Supérieure Polytechnique - Département de Génie Informatique — LIMSI
% Année       : 2025-2026
% Description : Fichier principal du cours. Orchestre tous les chapitres,
%               la bibliographie et les éléments de couverture.
% Compilateur : pdflatex + biber (ou lualatex recommandé)
%               → lualatex Cours && biber Cours && lualatex Cours && lualatex Cours
% ==============================================================================

\documentclass[12pt, a4paper, twoside, openright]{book}

% ==============================================================================
% CHARGEMENT DE LA FEUILLE DE STYLE
% ==============================================================================
\usepackage{Style_cours}

% ==============================================================================
% BIBLIOGRAPHIE
% ==============================================================================
\addbibresource{bibliographie.bib}

% ==============================================================================
% MÉTADONNÉES PDF
% ==============================================================================
\hypersetup{
    pdftitle  = {Audit des Systèmes d'Information},
    pdfauthor = {MINKA MI NGUIDJOÏ Thierry Emmanuel},
    pdfsubject= {Cours universitaire — Audit SI, Gouvernance, Sécurité},
    pdfkeywords={Audit SI, CISA, COBIT, ISO 27001, Contrôle interne,
                 Sécurité des systèmes d'information, Université de Yaoundé I}
}

% ==============================================================================
% DÉBUT DU DOCUMENT
% ==============================================================================
\begin{document}

% ------------------------------------------------------------------------------
% 1. PAGE DE TITRE (1ère DE COUVERTURE)
% ------------------------------------------------------------------------------
\begin{titlepage}
\thispagestyle{empty}
\pagecolor{CouleurPrimaire}

\begin{tikzpicture}[remember picture, overlay]
    % Bande dorée en bas de couverture
    \fill[CouleurAccent]
        (current page.south west) rectangle
        ([yshift=2.2cm] current page.south east);
    % Ligne décorative en haut
    \fill[CouleurAccent]
        ([yshift=-1.5cm] current page.north west) rectangle
        ([yshift=-1.8cm] current page.north east);
    % Élément graphique décoratif (carré tourné)
    \fill[white, opacity=0.04]
        ([xshift=8cm, yshift=-6cm] current page.north west)
        -- ++(5,0) -- ++(0,-5) -- ++(-5,0) -- cycle;
    \fill[white, opacity=0.06]
        ([xshift=9cm, yshift=-7cm] current page.north west)
        -- ++(6,0) -- ++(0,-6) -- ++(-6,0) -- cycle;
\end{tikzpicture}

\vspace*{1.5cm}

% Logo institutionnel (à décommenter si disponible)
% \begin{center}
%     \includegraphics[height=2.5cm]{figures/logo_uy1.png}
% \end{center}

\begin{center}
    {\color{white!80!gray}\small\scshape
        Université de Yaoundé~I~·~Département d'Informatique~·~LIMSI}\\[0.4ex]
    {\color{CouleurAccent}\rule{0.55\linewidth}{0.8pt}}
\end{center}

\vspace{1.5cm}

\begin{center}
    {\color{white!60!gray}\footnotesize\scshape
        cours universitaire}\\[1ex]
    {\color{white}\fontsize{34}{42}\selectfont\bfseries
        Audit des}\\[0.3ex]
    {\color{CouleurAccent}\fontsize{34}{42}\selectfont\bfseries
        Systèmes d'Information}\\[1.5ex]
    {\color{white!85!gray}\Large\itshape
        Fondements, Méthodes, Référentiels et Pratique}
\end{center}

\vspace{2cm}

\begin{center}
    {\color{white!70!gray}\small
        \ding{43}\quad Cas fil rouge : Entreprise LiZbiZ}\\[0.5ex]
    {\color{white!70!gray}\small
        \ding{44}\quad COSO~·~COBIT~·~ISO 27001~·~ISO 19011}\\[0.5ex]
    {\color{white!70!gray}\small
        \ding{228}\quad Audit fonctionnel~·~Audit de sécurité~·~Gouvernance IA}
\end{center}

\vfill

\begin{center}
    {\color{CouleurAccent}\rule{0.7\linewidth}{1.2pt}}\\[1.2ex]
    {\color{white}\large\bfseries MINKA MI NGUIDJOÏ Thierry Emmanuel}\\[0.4ex]
    {\color{white!80!gray}\small
        CISA~·~CISM~·~CGEIT~·~CRISC~·~CDPSE}\\[0.3ex]
    {\color{white!70!gray}\footnotesize\itshape
        Chercheur — LIMSI, Université de Yaoundé~I}\\[0.3ex]
    {\color{white!70!gray}\footnotesize\itshape
        Expert Judiciaire~—~Cadre supérieur, Contrôle Supérieur de l'État~—~Consultant International}\\[1.2ex]
    {\color{CouleurAccent}\rule{0.7\linewidth}{1.2pt}}
\end{center}

\vspace{0.6cm}

\begin{center}
    {\color{white!60!gray}\small Année académique 2024--2025}
\end{center}

\afterpage{\nopagecolor}
\end{titlepage}

% ------------------------------------------------------------------------------
% 2. PAGE DE DROITS / VERSO DE LA COUVERTURE
% ------------------------------------------------------------------------------
\clearpage
\thispagestyle{empty}
\vspace*{\fill}

\begin{center}
    \begin{minipage}{0.78\textwidth}
        \footnotesize\color{CouleurBordInfo}
        \begin{center}
            {\bfseries\color{CouleurPrimaire}
                Audit des Systèmes d'Information}\\[0.4ex]
            {\itshape Fondements, Méthodes, Référentiels et Pratique}
        \end{center}
        \vspace{0.8ex}
        \hrule
        \vspace{0.8ex}
        \textbf{Auteur :} MINKA MI NGUIDJOÏ Thierry Emmanuel\\
        \textbf{Institution :} Université de Yaoundé~I — LIMSI\\
        \textbf{Année :} 2024--2025\\[0.5ex]
        Ce document est un support de cours universitaire rédigé à des fins
        pédagogiques. Toute reproduction partielle ou totale à des fins commerciales
        est interdite sans l'autorisation écrite de l'auteur. Les références à des
        produits, organisations ou normes sont citées à titre pédagogique uniquement.\\[0.5ex]
        Les certifications CISA, CISM, CGEIT, CRISC et CDPSE sont des marques
        déposées de l'\textbf{ISACA} (\textit{Information Systems Audit and Control Association}).\\[0.5ex]
        \textbf{Contact :} Université de Yaoundé~I, Département d'Informatique,
        Yaoundé, Cameroun.\\[0.8ex]
        \hrule
        \vspace{0.5ex}
        {\centering\itshape
        «~Éduquer, c'est donner à l'autre les armes de sa propre lucidité.~»}
    \end{minipage}
\end{center}

\vspace*{\fill}

% ------------------------------------------------------------------------------
% 3. DÉDICACE
% ------------------------------------------------------------------------------
\clearpage
\thispagestyle{empty}
\vspace*{3cm}
\begin{flushright}
    \begin{minipage}{0.6\textwidth}
        \itshape\large\color{CouleurPrimaire}
        À ceux qui ont pris la peine de comprendre\\
        avant d'agir, et de vérifier avant de conclure.\\[1.5ex]
        \normalsize\normalfont\color{CouleurAccent}
        Et à tous les auditeurs en devenir\\
        qui liront ces pages avec l'exigence\\
        qu'ils apporteront demain à leur métier.
    \end{minipage}
\end{flushright}

% ------------------------------------------------------------------------------
% 4. REMERCIEMENTS
% ------------------------------------------------------------------------------
\clearpage
\thispagestyle{empty}
\chapter*{Remerciements}
\addcontentsline{toc}{chapter}{Remerciements}

L'élaboration de ce cours est le fruit d'un parcours intellectuel nourri par de nombreuses rencontres et expériences. L'auteur exprime sa reconnaissance au \textbf{Professeur Flavien Serge Mani Onana} et au \textbf{Professeur Thomas Djotio Ndié} pour leur encadrement scientifique rigoureux et leurs orientations éclairées dans la conduite des travaux de recherche.

Ses remerciements vont également à ses collègues des \textbf{Services du Contrôle Supérieur de l'État}, dont le pragmatisme institutionnel a confronté la théorie à la réalité du terrain, et à l'ensemble de l'équipe du projet \textbf{GIP Bénin}, dont les défis de gouvernance ont enrichi la substance de cet enseignement.

Enfin, ce cours n'aurait pas existé sans les étudiants qui, par leurs questions, leurs doutes et leur curiosité, en ont constamment éprouvé la clarté et la pertinence. Qu'ils en soient remerciés.

% ------------------------------------------------------------------------------
% 5. TABLE DES MATIÈRES
% ------------------------------------------------------------------------------
\clearpage
\tableofcontents

% ------------------------------------------------------------------------------
% 6. LISTE DES FIGURES (si nécessaire)
% ------------------------------------------------------------------------------
\clearpage
\listoffigures
\addcontentsline{toc}{chapter}{Liste des figures}

% ------------------------------------------------------------------------------
% 7. LISTE DES TABLEAUX
% ------------------------------------------------------------------------------
\listoftables
\addcontentsline{toc}{chapter}{Liste des tableaux}

% ------------------------------------------------------------------------------
% 8. LISTE DES CODES SOURCE
% ------------------------------------------------------------------------------
\lstlistoflistings
\addcontentsline{toc}{chapter}{Liste des codes source}

% ==============================================================================
% CORPS DU COURS
% ==============================================================================
\mainmatter

% --- Introduction Générale ---
\input{chapitres/0_Introduction}

% -------------------------------------------------------------------
% PARTIE I — Fondements Théoriques et Méthodologiques
% -------------------------------------------------------------------
\part{Fondements Théoriques et Méthodologiques}

% Chapitre 1 — Cas pratique fil rouge
\input{chapitres/01_cas_lizbiz}

% Chapitre 2 — Définitions, concepts et bases
\input{chapitres/02_fondamentaux}

% Chapitre 3 — Démarche générale d'audit SI
\input{chapitres/03_demarche}

% -------------------------------------------------------------------
% PARTIE II — Référentiels et Typologies d'Audit
% -------------------------------------------------------------------
\part{Référentiels et Typologies d'Audit}

% Chapitre 4 — Normes et référentiels usuels (COSO, COBIT, ISO 27000)
\input{chapitres/04_referentiels}

% Chapitre 5 — Principaux types d'audit SI
\input{chapitres/05_types_audits}

% -------------------------------------------------------------------
% PARTIE III — Enjeux Contemporains et Évaluation
% -------------------------------------------------------------------
\part{Enjeux Contemporains et Évaluation}

% Chapitre 6 — Sujets connexes : IA et Fraude
\input{chapitres/06_sujets_connexes}

% Chapitre 7 — Évaluation finale
\input{chapitres/07_evaluation_finale}

% -------------------------------------------------------------------
% TRAVAUX PRATIQUES (Laboratoires)
% -------------------------------------------------------------------
\part{Travaux Pratiques Immersifs}

\input{chapitres/lab1_audit_applicatif}
\input{chapitres/lab2_infrastructure}

% --- Conclusion Générale ---
\input{chapitres/00_Conclusion}

% ==============================================================================
% MATIÈRE POSTLIMINAIRE
% ==============================================================================
\backmatter

% --- Bibliographie ---
\clearpage
\printbibliography[
    heading=bibintoc,
    title={Bibliographie et Références}
]

% --- À Propos de l'Auteur ---
\clearpage
\input{chapitres/Apropos_Auteur}

% --- Index ---
\clearpage
\printindex
\addcontentsline{toc}{chapter}{Index}

% ==============================================================================
% 4ÈME DE COUVERTURE
% ==============================================================================
\clearpage
\thispagestyle{empty}
\pagecolor{CouleurPrimaire}

\begin{tikzpicture}[remember picture, overlay]
    \fill[CouleurAccent]
        (current page.south west) rectangle
        ([yshift=1.8cm] current page.south east);
    \fill[CouleurAccent]
        ([yshift=-1.4cm] current page.north west) rectangle
        ([yshift=-1.7cm] current page.north east);
\end{tikzpicture}

\vspace*{2.5cm}

\begin{center}
    {\color{CouleurAccent}\rule{0.6\linewidth}{1pt}}\\[1.5ex]
    {\color{white}\Large\bfseries À propos de ce cours}\\[1.5ex]
    {\color{CouleurAccent}\rule{0.6\linewidth}{1pt}}
\end{center}

\vspace{1cm}

\begin{center}
\begin{minipage}{0.75\textwidth}
\color{white!85!gray}\setstretch{1.4}
Ce cours propose une approche intégrée et rigoureuse de l'audit des systèmes d'information, articulant les fondements théoriques du contrôle interne, les référentiels internationaux reconnus (COSO, COBIT, ISO 27001) et une pratique immersive à travers le cas fil rouge de l'entreprise LiZbiZ.

\medskip

Il s'adresse aux étudiants en informatique, gestion et sécurité des SI souhaitant acquérir les compétences opérationnelles et conceptuelles requises pour conduire des missions d'audit avec professionnalisme, en conformité avec les standards de l'ISACA et de l'IIA.

\medskip

Les thèmes abordés couvrent la démarche d'audit, l'analyse des risques, les tests fonctionnels et techniques, la gouvernance de l'intelligence artificielle, la détection de la fraude et la rédaction du rapport final.
\end{minipage}
\end{center}

\vfill

\begin{center}
    {\color{CouleurAccent}\rule{0.5\linewidth}{0.8pt}}\\[1ex]
    {\color{white}\bfseries MINKA MI NGUIDJOÏ Thierry Emmanuel}\\[0.3ex]
    {\color{white!75!gray}\footnotesize
        CISA · CISM · CGEIT · CRISC · CDPSE}\\[0.3ex]
    {\color{white!65!gray}\footnotesize\itshape
        Chercheur — LIMSI, Université de Yaoundé~I · Cameroun}\\[1ex]
    {\color{CouleurAccent}\rule{0.5\linewidth}{0.8pt}}
\end{center}

\vspace{1cm}

\afterpage{\nopagecolor}

\end{document}
% ==============================================================================
% Fin de Cours.tex
% ==============================================================================
